\section{Ускорение на графическом процессоре средствами SYCL/AdaptiveCpp}

\subsection{Слой compute/: API и архитектура}

Слой \texttt{compute/} инкапсулирует все обращения к вычислительному стеку графического процессора. Физическое ядро взаимодействует только с этим
слоем и никогда не использует SYCL напрямую. Это позволяет сменить технологию вычислений, не затрагивая физическую 
логику.

Основные компоненты слоя:

\paragraph{\texttt{Device}.} Представляет вычислительное устройство. Устройство выбирается статическими 
методами \texttt{default\_gpu()} и \texttt{default\_cpu()}; через него создаётся \texttt{Stream}.

\paragraph{\texttt{Stream}.} Обёртка над \texttt{sycl::queue} в режиме \textit{in-order}. Все операции, 
поставленные в один \texttt{Stream}, выполняются строго последовательно. Метод \texttt{wait()} блокирует центральный процессор
до завершения всех поставленных задач.

\paragraph{\texttt{Buffer<T>}.} Буфер в USM-памяти устройства. Предоставляет методы \texttt{upload} и 
\texttt{download} для асинхронной передачи данных ЦП$\leftrightarrow$ГП; при уничтожении буфера память 
освобождается автоматически.

\paragraph{\texttt{parallel\_for(s, n, f)}.}
Запускает вычислительное ядро: лямбда \texttt{f(size\_t i)} вызывается
для каждого $i \in [0, n)$ параллельно.
Ядро должно быть тривиально копируемым; внутри допустимы только device-safe операции.

\paragraph{\texttt{reduce\_sum}, \texttt{sort\_by\_key}.}
Детерминированные примитивы: редукция возвращает одно и то же значение
при одинаковых входных данных на одном устройстве.

\subsection{Жизненный цикл данных}

Все горячие данные живут в памяти устройства между шагами симуляции: ни позиции, ни скорости тел не 
переносятся на центральный процессор при каждом шаге. Единственная синхронизация за кадр между центральным и графическим процессором~--- чтение счётчика 
кандидатных пар после широкой фазы (одно 4-байтное чтение).

\begin{enumerate}
  \item При инициализации сцены данные копируются ЦП$\to$ГП один раз (\texttt{upload}).
  \item На каждом шаге симуляции вычислительные ядра модифицируют данные in-place.
  \item Для сохранения траектории (файл \texttt{.trajectory}) данные
    выгружаются ГП$\to$ЦП один раз на два кадра; это единственная загрузка шины.
  \item По завершении симуляции финальное состояние выгружается для
    вычисления контрольного хеша.
\end{enumerate}

\subsection{Целевые платформы AdaptiveCpp}

AdaptiveCpp компилирует одни и те же C++20-исходники для нескольких вычислительных платформ:

\begin{table}[h]
\centering
\caption{Бэкенды AdaptiveCpp и применяемые технологии}
\label{tab:backends}
\begin{tabular}{lll}
\toprule
\textbf{Пресет CMake} & \textbf{Бэкенд} & \textbf{Оборудование} \\
\midrule
\texttt{cuda}    & NVIDIA CUDA PTX  & NVIDIA RTX / Tesla / A-серия \\
\texttt{hip}     & AMD HIP/ROCm     & AMD RDNA / CDNA \\
\texttt{omp}     & OpenMP (ЦП)      & Любой многоядерный процессор \\
\bottomrule
\end{tabular}
\end{table}

Выбор бэкенда влияет только на флаги компиляции, но не на исходный код движка.

\subsection{Детерминизм на уровне примитивов}

Стандартная атомарная операция \texttt{atomicAdd} в CUDA/HIP не детерминирована: результат зависит от порядка 
прибытия нитей. Для детерминированных редукций в \texttt{compute/} реализован примитив \texttt{reduce\_sum}, 
использующий дерево редукции с фиксированным порядком операций. Аналогично, широкая фаза после построения 
дерева выполняет детерминированную сортировку кандидатных пар, прежде чем передать их в узкую фазу. Благодаря 
этим мерам одна и та же сцена даёт побитово идентичный результат при каждом запуске на одном устройстве и 
одной сборке.

\subsection{Выводы по разделу}

Слой \texttt{compute/} обеспечивает три ключевых свойства системы: изоляцию физической логики от API графического процессора (что
позволяет тестировать её на центральном процессоре), поддержку нескольких платформ с графическими процессорами из единого исходного кода, и детерминизм
вычислений через специализированные примитивы.
